\begin{lemma}[Well-foundedness of copied pair descent]
The relation $\mathsf{PowerGameDescent}$ on pairs of
$\mathsf{PowerQ}(Q)$-presentations is well-founded.
\end{lemma}
\begin{proof}
Write $R$ for the child-to-parent relation from \noderef{PowerChild}; its
orientation is that $R(a',a)$ makes $a'$ a predecessor of $a$.  By
\noderef{PowerChildWellFounded}, $R$ is well-founded.  We prove the desired
statement by well-founded induction on the first coordinate, with motive
\[
  \mathcal P(a)\;:\!\iff\;\text{for every }b,
  \mathsf{Acc}(\mathsf{PowerGameDescent},(a,b)).
\]
Thus fix $a$ and assume as outer induction hypothesis that, whenever
$R(a',a)$, every pair $(a',b')$ is accessible for
$\mathsf{PowerGameDescent}$.  It remains to prove $\mathcal P(a)$.  Fix an
arbitrary second coordinate $b$ and perform well-founded induction on $b$
with respect to the same relation $R$.  Its inner induction hypothesis says
that whenever $R(b',b)$, the pair $(a,b')$ is accessible.

We show that $(a,b)$ is accessible.  Let $(a',b')$ be a predecessor of it
for $\mathsf{PowerGameDescent}$.  By the definition in
\noderef{PowerGameDescent}, this means either
\[
  R(a',a),
\]
or
\[
  a'=a\quad\text{and}\quad R(b',b).
\]
In the first case, the outer induction hypothesis gives accessibility of
$(a',b')$, with no restriction on the second coordinate.  In the second
case, the inner induction hypothesis gives accessibility of $(a,b')$, and
the equality $a'=a$ transports this to accessibility of $(a',b')$.
These are all predecessors because the displayed two cases are precisely
the disjunction defining \noderef{PowerGameDescent}.  The accessibility
constructor therefore proves that $(a,b)$ is accessible.  The inner
induction proves this for every $b$, and the outer induction then proves it
for every $a$, which is exactly well-foundedness of
$\mathsf{PowerGameDescent}$.
\end{proof}
